% Nouveau EGA V (cote 25): archive cote 25, site page 154 (PDF page 155).
% typescript. Draft compared: none.
% AI-transcribed by Claude Opus 5.5 (claude-opus-5-5, adaptive thinking, effort high), Message Batches API, batch ega25b (msgbatch_01A7n5YuXubUPNK1zsepi8RG), 2026-11-13; original AI transcription, subsequently checked against both supplied scan parts; see PAGES_CHECKED.tsv and CHECK_DECISIONS.tsv.
\scanpage{25/154}
\mspage{\hand{15}}

Pour caractériser par voie différentielle  les  zéros singuliers ordinair\unsure{es}
relativement à \over{S}{Y},  on note qu'il faut simplement exprimer  les conditions
de la proposition 9  pour  l'homomorphisme \xout{naturel} déduit de u
\[ \xout{d^{\hand{\infty}}\illegible{1}E} \qquad P^{\hand{\infty}}_{X/\over{S}{Y}}(\underline{E})_{V^1}\hand{\to}(P^{\hand{\infty}}_{X/\over{S}{Y}})_{V^1} \quad , \]
qui se réduit d'ailleurs à une condition  ``d'ordre 2'' i.e. portant seule-%
ment sur l'homomorphisme correspondant
\[ P^{2}_{X/\over{S}{Y}}(\underline{E})_{V^1}\hand{\to}(P^{2}_{X/\over{S}{Y}})_{V^1} \ . \]
Il faut exprimer en premier lieu  que  le noyau \inserted[typed]{$\underline{K}$} de
\[ (\underline{x}\ ) \qquad \underline{E}_{V^1} \hand{\longrightarrow} (\underline{\hand{\Omega}}^{1}_{X/\over{S}{Y}})_{V^1} \]
en x est inversible et facteur direct, i.e. que l'homomorphisme en\unsure{v|i}\unsure{i|v}s\unsure{a|g}gé
(qui en vertu de  $x\hand{\in} V^1$ est  de rang  m-1 en x) est exactement de rang
\margin{(} m-1 en x  (NB \uline{le rang} d'un  homomorphisme de \unsure{m}odules loc libres en un
point est par définition le rang de 'lhomomorphisme de vect\unsure{o|r}\unsure{r|i}\unsure{i|r}\unsure{e|d}ls sur k(x)
correspondant), c e qui s'exprime en considérant
\[ (\hand{\bigwedge}^{m-1} \underline{E})_{V^1}\hand{\to}(\underline{\hand{\Omega}}^{m-1}_{X/\over{S}{Y}})_{V^1} \quad , \]
le  schéma  V'' des zéros de cet homomorphisme  ( $V''\hand{\subset} V^1$ )  et en écrivant
simplement $x\hand{\notin}$ V'' .  La définition de V'' assure que  K $|V^1-V''$  est un
Module inversible,  et que  le conoyau  $\underline{W}$  de ($\underline{x}$ )  est tel que
$\underline{W}\,|V^1-V''$  est  localement libre de rang égal à  \over{n}{d}-(m-1)  (où \over{n}{d} est la
dimension relative de X sur S, égale au rang de $\underline{\hand{\Omega}}^{1}_{X/\over{S}{Y}}$).  Appliquant la
construction expliquée avec la prop.9\hand{$^{*}$}, on trouve \xout{maintenant} un homomor-%
phisme
\[ M\ : \qquad \underline{K} \hand{\longrightarrow} \underline{\mathrm{Sym}}\xout{m}^{2}(\underline{W}) \]
et il reste à exprimer que ce dernier est ``non dégénéré'' en x. Pour ceci,
introduisons le  discriminant
\[ D\ \in\ \hand{\Gamma}\,(V^1-V'',\ \underline{K}^{-n+m\over{\unsure{+|-}}{\unsure{-|+}}1}\otimes \text{dét}(\underline{W})^{2}\ ) \]
(où les exposants désignent des puissances tensorielles de Modu\unsure{l|b}es

\mspage{\hand{154}}
